\documentclass[11pt]{article}
\usepackage[utf8]{inputenc}
\usepackage[T1]{fontenc}
\usepackage{geometry}
\usepackage{titlesec}
\usepackage{enumitem}
\usepackage{hyperref}

\geometry{margin=1in}
\titleformat{\section}{\large\bfseries\uppercase}{}{0em}{}
\titleformat{\subsection}{\bfseries}{}{0em}{}

\title{Fate's Edge: Streamlined Lorebook\\(Areas Only — No Akilan or Ostrilan)}
\author{}
\date{}

\begin{document}

\maketitle

\section*{Geographic Overview (No Real‐World Analogues)}
The known world is a mosaic of seas, mountains, forests, and steppes.  
The \textbf{Amaranthine Sea} washes the southern coasts – warm, trade‑rich, dotted with islands and peninsulas.  
The \textbf{Dolmis} (eastern sea) is colder, rougher; its shores are a patchwork of clans, city‑states, and fallen imperial provinces.  
North of the Aelerian Mountains (a high, impassable spine) lie the \textbf{Mistlands}, the \textbf{Direwood}, and the \textbf{Violet Steppe} – a foggy breadbasket, a hungry forest, and a wind‑scoured grassland.  
Beyond the steppe, the \textbf{Aberderrin Sea} mirrors the north with iron skies and brackish waves.  
The great forests of the \textbf{Valewood} and the gnomish shores of \textbf{Averossa} flank the eastern Dolmis, while the halfling downs of \textbf{Amedell} nestle south of them.  
The southern continent, \textbf{Akilan}, lies beyond the Titan’s Throat strait – a land of monsoons, walking cities, and ancient empires.  

All descriptions below avoid real‐world names; geography is expressed through climate, terrain, and cultural logic.

\section*{1. Utar (The Utaran Imperium \& Successor Northlands)}

\subsection*{Overview}
Once a vast empire that ruled from the Amaranthine to the northern foothills, Utar fractured into quarrelling successor states. Now the north is a land of crumbling marble roads, disputed charters, and the bitter memory of order.

\subsection*{Key Regions}
\begin{itemize}[leftmargin=*]
    \item \textbf{Ecktoria} – The imperial heart, now a relic‑state. White forums, gold‑sealed coin‑houses, and a populace that recites laws as prayers.
    \item \textbf{Acasia} – A broken march of petty warlords and the cosmopolitan city of \textbf{Silkstrand}, where bridges are counted in debts and dyewater never washes out.
    \item \textbf{Vhasia} – The fractured sun. Duchies that claim a dead crown, where every oath is sworn twice and broken once. Parlement clerks duel with citations.
    \item \textbf{Viterra} – The last true kingdom. Hedged fields, parish stones, and a warrior queen who judges by precedent and longsword.
    \item \textbf{Thepyrgos} – A city of a thousand stairs climbing a cliff. Bell‑law, synod debates, and high‑elf archivists who count the breaths of scholars.
\end{itemize}

\section*{2. Kahfagia — The Empire of Wakes and Storm-Flags}

\subsection*{Overview}
A maritime oligarchy straddling the Titan’s Throat, the narrow strait between the Amaranthine and the southern seas. Kahfagia lives by trade, tides, and the lantern‑law that shifts with the salt wind.

\subsection*{Key Features}
\begin{itemize}[leftmargin=*]
    \item \textbf{Kassamira} – A city of white lime and black rigging; its harbor mouth guarded by a boom chain blessed by kraken‑priests.
    \item \textbf{Stormspire} – A fortified lighthouse‑fortress where admirals plot privateer campaigns against Ecktorian grain convoys.
    \item Lantern‑law: jurisdiction flows with the tide – at high water, the Admiralty rules; at ebb, the merchant courts.
    \item The sea is both highway and graveyard. Whales beached on sacrifice altars are omens of new trade routes or new wars.
\end{itemize}

\section*{3. Theona — The Marsh Crown}

\subsection*{Overview}
Three green isles in the Dolmis, pledged to Viterra in name only. Theona is a place of peat smoke, standing stones, and the eerie \textit{No Ninth} taboo.

\subsection*{Key Features}
\begin{itemize}[leftmargin=*]
    \item Bogs that swallow sound; hedged fields where scarecrows watch the lane.
    \item \textbf{The Ninth Taboo:} Never pour the ninth cup, never count the ninth step, never speak the ninth name. Omissions are courtesies to the Hollow.
    \item Coracle fleets that slip through the reed‑mazes, their crews humming old hymns that skip the forbidden beat.
    \item The Moot Hill, where feuds are settled with a knife and a story, and where the wind always tastes of rain.
\end{itemize}

\section*{4. The Mistlands — Bells, Salt, and Breath}

\subsection*{Overview}
A fog‑drenched breadbasket north of the Aelerian Mountains, ruled by dwarven protectorates. The Mistlands feed the lowlands, but the mist demands its toll in memory and breath.

\subsection*{Key Features}
\begin{itemize}[leftmargin=*]
    \item Bell‑line levees: iron posts strung with bronze bells that ring at dusk, keeping the Direwood’s breath at bay.
    \item Reed‑maze causeways where a wrong turn means a year of wading through whispers.
    \item The \textbf{Weeping Gate} – a western palisade of bone‑white wood, wired with bells that cry when the forest exhales.
    \item Salt pans worked by monks who bless each grain; a handful of ward‑salt can buy a night’s safe passage.
\end{itemize}

\section*{5. Valewood — The Forest That Remembers You Wrong}

\subsection*{Overview}
A sentient forest east of the Aberderrin Sea, where the ruins of a fallen elven empire phase in and out of reality. The Valewood does not forgive trespass; it re‑members you into its logic.

\subsection*{Key Features}
\begin{itemize}[leftmargin=*]
    \item Phasing ruins: star‑roads that appear only under certain stars, amphitheatres that remember speeches not yet given.
    \item Trees with faces that watch, roots that write contracts in the dirt.
    \item The \textbf{Lethai‑ar} (wood elves) who keep the old courts – the Hazel Queen and the Alder King – and whose masks never come off in the presence of iron.
    \item \textbf{Green Neighbors} – fae or fey‑touched beasts that offer shortcuts in exchange for a name, a memory, or a year of silence.
\end{itemize}

\section*{6. Ykrul — Storm on the Steppe}

\subsection*{Overview}
Nomadic clans of horse‑riders and raiders who follow the grass and the omens. Their law is the \textbf{Kon'reh} – sacred geometry of routes, exits, and blood‑price.

\subsection*{Key Features}
\begin{itemize}[leftmargin=*]
    \item Seasonal migrations from the Violet Steppe to the Aberderrin shore. Encampments are rings of felt wagons, cookfires outward.
    \item Shamans who read the wind and the marrow; the \textbf{Faith of the Open Sky} has no temples, only horizont.
    \item The Kurultai – a council of clan mothers where feuds are settled by horse‑races and bone‐lots.
    \item Dragon boats dragged over portages, launched into the Yloka River to raid the Dolmis coast.
\end{itemize}

\section*{7. Zakov — Salt \& Serpent}

\subsection*{Overview}
A pirate haven on a jagged island in the Dolmis. Zakov is ruled by the Seven Guilds, each controlling a dock, a trade, or a secret. Law is a suggestion; the tide is the only judge.

\subsection*{Key Features}
\begin{itemize}[leftmargin=*]
    \item The \textbf{Crimson Docks} – contracts signed in blood, enforced by the Syndicate’s collectors.
    \item \textbf{Serpent’s Spine} – a reef that shifts with the moon; only native pilots know the safe cuts.
    \item The Salt Prince, a masked figure who auctions letters of marque while his shadow votes against him.
    \item Sunken temples in the harbour, where drowned debtors whisper secrets to those brave enough to listen.
\end{itemize}

\section*{8. Ubral — The Stone Between Spears}

\subsection*{Overview}
Highland clans, cairns, and iron oaths. Ubral is a land of shepherds and reivers, caught between Viterran order and Vhasian chaos. Every hill has a name, and every name has a feud.

\subsection*{Key Features}
\begin{itemize}[leftmargin=*]
    \item \textbf{Khaz‑Vurim} – a dwarven mountain‑hold that carves its tolls into the pass; the stone remembers every unpaid debt.
    \item Guest‑right tokens: a bronze ring that guarantees three nights’ shelter and a sword at your back.
    \item \textbf{The Pass of Ashes} – an old volcano road where the wind tastes of cinders and ghosts of legionaries still drill.
    \item Cairn courts: disputes settled atop burial mounds, the ancestors as silent jurors.
\end{itemize}

\section*{9. Linn — Skerries \& Storm-Oaths}

\subsection*{Overview}
Southernmost Linnic tribes, raiders of the Dolmis coast and river‑landers of the Yloka. They swear oaths on whale‑bone and keep their sagas in braided hair.

\subsection*{Key Features}
\begin{itemize}[leftmargin=*]
    \item Fjords, skerries, and longship sheds that smell of tar and cold iron.
    \item The \textbf{Thing‑holm} – a flat islet where law is spoken, feuds are priced in silver, and the sea listens.
    \item The Sea‑Queen’s court – a driftwood throne, her crown a circlet of wolf teeth.
    \item \textbf{Whale‑road} – the open sea route to the Aberderrin, dangerous with ice and white squalls.
\end{itemize}

\section*{10. Aelinnel — Stone, Bough, and Bright Things (Averossa)}

\subsection*{Overview}
The gnomish homeland of Aevrossa, a rugged coast of granite, hawthorn, and tide‑cut stairs. The Aelinnel count in primes, favour copper over iron, and keep two ledgers for every promise.

\subsection*{Key Features}
\begin{itemize}[leftmargin=*]
    \item Dolmen stairs and basalt organ cliffs that “breathe” in prime intervals.
    \item Hedge‑witches who tie red threads to bind oaths – or to curse the oath‑breaker.
    \item The \textbf{Green Gate} – a sea‑arch that opens at certain tides to a forest not on any map.
    \item Thorn Court rings where the Lady of Thorns judges with petals that cut like knives.
\end{itemize}

\section*{11. Aelaerem — Hearth \& Hollow (Amedell)}

\subsection*{Overview}
The halfling downs of Amedell – orchards, thatched cottages, and lanes that lead to standing stones. The Aelaerem live by hospitality, lantern‑law, and the careful omission of the ninth cup.

\subsection*{Key Features}
\begin{itemize}[leftmargin=*]
    \item Moot Oaks with iron nails hammered deep; beneath them, the Apple‑Matron pours cider and settles disputes.
    \item Hedge‑witches who read omens in the flight of watch‑geese.
    \item Mummers’ processions where masks become faces, and after dark, the Hollow walks.
    \item The \textbf{Pale Shepherd} – a spirit of thresholds who guides lost children home in exchange for a story.
\end{itemize}

\section*{12. Aeler — Crowns \& Under-Vaults}

\subsection*{Overview}
Dwarven mountain holds beneath the Aelerian peaks. Aeler is a realm of breath‑ledgers, keystone rights, and the patient weight of stone.

\subsection*{Key Features}
\begin{itemize}[leftmargin=*]
    \item Vaultmouth gates that test air quality before admitting strangers.
    \item Under‑roads lit by brass lamps; every light is counted, every breath a debt.
    \item Spirit Shield Warriors – ancestor‑venerating protectors whose masques are heirlooms of the fallen.
    \item The High King Beneath the Peaks, who has not seen the sun in three hundred years – but still speaks through bell‑codes.
\end{itemize}

\section*{13. Black Banners — Condotta \& Crowns}

\subsection*{Overview}
Mercenary lands and war camps that shift with the season. The Black Banners are not a nation but a contract; their loyalty is ink, and ink fades.

\subsection*{Key Features}
\begin{itemize}[leftmargin=*]
    \item The \textbf{Singing Wastes} – a battlefield where fallen weapons hum in the wind, each note a dying oath.
    \item Condotta brokers who sell companies by the month; their seals change colour with every payday.
    \item The Bannerless One – a legendary commander who commands loyalty without flag or coin; some say they are a ghost, others a prophecy.
    \item Bone Fields where Ykrul elders parley among skulls, reckoning debts that outlive the flesh.
\end{itemize}

\section*{14. Vilikari — Laurels \& Longhouses}

\subsection*{Overview}
A federated frontier where Utaran law meets barbarian custom. The Vilikari keep two laws: the wolf (blood‑price) and the eagle (imperial statute). Both are written on the same stone.

\subsection*{Key Features}
\begin{itemize}[leftmargin=*]
    \item The \textbf{Foedus Stone} – a crossroads monolith where treaties are renewed each spring with oaths and horse‑sacrifice.
    \item Twin courts: the mallus (oak‑ring, clan elders) and the basilica (marble bench, imperial clerks). Their verdicts rarely agree.
    \item March towns of timber and tile, where a man can be a barbarian in the morning and a citizen by noon.
    \item The Queen of the Marches, who wears both a laurel wreath and a wolf‑tail cloak.
\end{itemize}

\section*{15. The Wilds — Roads, Ruins, and Weather}

\subsection*{Overview}
Untamed lands that change by biome: forest, desert, tundra, swamp, coast. The Wilds are a reskin palette – but always, they remember.

\subsection*{Key Features}
\begin{itemize}[leftmargin=*]
    \item Crossing points: fords, ice‑spans, stepping logs, dune saddles – each with its own price.
    \item Shelter hollows: a cave, a root cellar, a fallen giant’s ribcage. Soot on the ceiling means someone survived the night.
    \item Forager children who trade route‑songs for a scrap of bread; their maps are drawn in charcoal and always accurate until sunset.
    \item Weather as a character: glass‑clear haze, white squall, dust devil, ground mist. The sky has moods, and it shares them freely.
\end{itemize}

\section*{16. Tulkani — Road-Kin of the Ember Line}

\subsection*{Overview}
Nomadic clans of painted wagons, braided oaths, and ember stories. Once they had a homeland; now the road is their hearth.

\subsection*{Key Features}
\begin{itemize}[leftmargin=*]
    \item Wagon rings drawn at dusk – the circle is sanctuary, the gap is an invitation to the Hollow.
    \item Fire‑cults that sing the names of the dead, one verse per log.
    \item Kuva of the Hearth‑Road: a spirit of journey and farewell, worshipped in the glow of coals, not temples.
    \item The Family of the Raven Road – a sept known for carrying whispers between armies; they never betray a secret, but they always charge interest.
\end{itemize}

\section*{17. Ikari — First Plough, First Oath}

\subsection*{Overview}
Native tillers and smiths of the northern continent, older than the Utaran conquest. The Ikari keep to their ringforts and ancestor fires, waiting for the empire’s shadow to lift.

\subsection*{Key Features}
\begin{itemize}[leftmargin=*]
    \item Tribes like the Kreki (fishermen of the Aberderrin) and the Smeinnoii (smiths of the eastern hills).
    \item Ancestral fires kept burning for generations; to let one go out is to invite a curse.
    \item Law‑keepers who recite the \textbf{Ondriti code} – a cycle of seasonal judgments carved into antler.
    \item Seasonal raids on lowland granaries, conducted with ritualised restraint (blood‑price is always offered).
\end{itemize}

\section*{18. Midh Adkaz — Where War Became a Market}

\subsection*{Overview}
A frontier camp that grew into a crossroads city. Midh Adkaz is a place where oaths are traded like grain and where yesterday’s enemy is today’s trading partner.

\subsection*{Key Features}
\begin{itemize}[leftmargin=*]
    \item The \textbf{Red Ditch} – a dry moat where duels are fought over disputed cargo; the loser pays, the winner pays the harbourmaster.
    \item The \textbf{Stakefield} fairground: merchants from five nations pitch tents, and the Six Hands council arbitrates with weighted dice.
    \item The Boar Gate and River Gate – two entrances that never lock, because every guard captains a competing syndicate.
\end{itemize}

\section*{19. Haayr Peninsula — Anvils Between Two Seas}

\subsection*{Overview}
Mountain tongues and broken coasts that separate the Amaranthine from the Dolmis. The Haayr Peninsula is a strategic chokepoint, and its passes are paved with the bones of armies.

\subsection*{Key Features}
\begin{itemize}[leftmargin=*]
    \item The Spine of Haayr – a limestone ridge with only three navigable passes, each controlled by a different fortress.
    \item \textbf{Khar‑Myra} – a port city carved into a sea cliff; its lower docks are underwater at high tide.
    \item Theressos – a temple‑city built into extinct volcanoes, where oracles breathe ash and foretell grain prices.
    \item The Hook Road, a switchback path where merchants hire acrobats to spot avalanches.
\end{itemize}

\section*{20. Dhahara — Monsoon of Empires}

\subsection*{Overview}
A land of oases, monsoons, and marching armies – a frontier where caravans from the steppe meet fleets from the Amaranthine.

\subsection*{Key Features}
\begin{itemize}[leftmargin=*]
    \item The \textbf{Himdal Marches} – terraced hills that catch the rain; every drop is tithed to the canal‑mothers.
    \item The \textbf{Jade Oases} – springs that change colour with the season, sacred to the Spice Lords.
    \item \textbf{Sarvash} – a city of suspended gardens, where nobles ride elephants through silk‑hung streets.
    \item The Monsoon Bells – bronze chimes hung in mountain passes; their tone tells merchants how many days until the roads drown.
\end{itemize}

\section*{21. Oshiira — The Ledger Empire}

\subsection*{Overview}
A confederation of canals and ledgers, built on the bones of a fallen river‑kingdom. Oshiira measures wealth in bushels, justice in weights, and honour in interest.

\subsection*{Key Features}
\begin{itemize}[leftmargin=*]
    \item Canal webs that criss‑cross the Sahel grasslands; every lock has a clerk, every clerk has a stamp.
    \item The \textbf{Spirit of the Long Sorrow} – a low, constant drone from underground aqueducts; it is said to be the breath of a drowned god.
    \item Mbaro, the capital – a city of weigh‑houses and ledgers; its skyline is flat because no building may exceed the Granary Tower.
    \item The Mbari‑style Senate, where votes are cast with stamped receipts and filibusters are drowned in ink.
\end{itemize}

\section*{22. Sekogo — Where the Roads Meet the Tide}

\subsection*{Overview}
A crossroads of river and sea, where the Crimson Basin drains into the southern ocean. Sekogo is a land of grove‑masters, lagoon wardens, and the Tide Ledger.

\subsection*{Key Features}
\begin{itemize}[leftmargin=*]
    \item The \textbf{Mbaro Quays} – floating markets tethered to mangrove roots; prices are sung, not shouted.
    \item \textbf{Brasswater Row} – a canal of patinated bronze gutters where alchemists brew tide‑charms.
    \item Quay Syndicates that control the spice trade; their symbols are knots, not flags.
    \item The Circle of Unkwa – a druidic council that reads the future in the trails of ghost crabs.
\end{itemize}

\section*{23. Taharka — Monsoon Crown, Terrace Throne}

\subsection*{Overview}
A highland kingdom of canals, convoys, and coins. Taharka is a realm of water and stone, where every drop is measured and every stone is mortgaged.

\subsection*{Key Features}
\begin{itemize}[leftmargin=*]
    \item The \textbf{Mkusaro Highlands} – terraced farms that spiral up the mountains; each terrace has its own microclimate.
    \item The \textbf{Canal Collegium} – a guild of hydraulic engineers who wear copper rings for every lock they’ve built.
    \item \textbf{Siatwe}, the capital – a city of spiraling bazaars and stepped water‑gardens; its mint produces coins that weigh truth.
    \item Meltwater courts: judges sit on stone chairs fed by canals; the flow of water determines the speed of justice.
\end{itemize}

\section*{24. Ameria — Between Bay and Throat}

\subsection*{Overview}
A divided realm between Kahfagia and the Titan’s Throat, perpetually neutral by treaty and perpetually raided by privateers. Ameria is a buffer of royal forms and empty treasuries.

\subsection*{Key Features}
\begin{itemize}[leftmargin=*]
    \item The \textbf{Shoreless Bay} – a broad, shallow sound where grain barges anchor; the mud is said to swallow secrets.
    \item Throatward Ports – fortified harbours that pay tribute to both Kahfagia and the southern kingdoms.
    \item \textbf{Khol‑Amar} – a hill‑fort of red sandstone, seat of the Neutral Regency.
    \item Cape Verdan, where lighthouse‑keepers burn green flares to confuse corsairs.
\end{itemize}

\section*{25. Ngomebe — Stone Men of the Moving Cities}

\subsection*{Overview}
Ironworkers and wall‑builders of Akilan, famous for their walking cities – colossal stone structures on rollers, dragged across the plateau by oxen and ancestral will.

\subsection*{Key Features}
\begin{itemize}[leftmargin=*]
    \item \textbf{Mkusaro ridges} – ancient quarry‑sites where the stone still sings when tapped at dawn.
    \item The Ekale Spurs – foothills dotted with half‑finished cities, abandoned when the keystone cracked.
    \item \textbf{Khazembo} – a walking city that has never stopped in living memory; its rollers are polished to mirrors.
    \item Mason‑Kin Houses – guilds that share a single underworld with Aeler engineers; they speak in hammer‑codes and ledger‑rhymes.
\end{itemize}

\section*{26. Ashaan — Gem of the Sea, Shadow on the River}

\subsection*{Overview}
A fallen slaver‑empire, now a magicocracy ruled by the Veiled Sisters. Ashaan is beautiful, decadent, and rotting from within; its spirit‑bound economy runs on stolen breath.

\subsection*{Key Features}
\begin{itemize}[leftmargin=*]
    \item The \textbf{Veiled, Helmed, and Masked Sisters} – three orders of sorcerer‑nobles who control the courts, the docks, and the nightmares.
    \item The Esoti – magistrates who rewrite ledgers with invisible ink; a name erased is a life forgotten.
    \item \textbf{Galanina} – a port city of black marble and brass cages, where the air hums with bound spirits.
    \item The Black Hand – secret police who wear no uniform; they could be your neighbour, your clerk, your lover.  
    \item The \textbf{Breath‑less} – a resistance network of escaped slaves and heretical mages, who whisper Ikasha’s name in root cellars.
\end{itemize}

\section*{27. Sihai — The Central Kingdom, The Ordered Land}

\subsection*{Overview}
An immense, ancient empire of rigid hierarchy, bureaucratic exams, and the Mandate of Heaven. Sihai is a land where mountains are measured in poems and rebellions are filed as paperwork.

\subsection*{Key Features}
\begin{itemize}[leftmargin=*]
    \item The \textbf{Sihon River Basin} – a fertile plain criss‑crossed by canals, guarded by garrison towns on stilts.
    \item The \textbf{Himadri Mountains} – the northern border, where earth‑shakers (spirit cultivators) train in hidden pavilions.
    \item The Bureaucracy – a vast corps of scholar‑officials who compete in poetry contests; the winner governs, the loser cleans ink‑stones.
    \item Warrior Monks who guard the passes and serve no lord – they believe the emperor is a temporary inconvenience.
    \item The \textbf{Hintara Ocean Coast} – a string of pearl‑fishing villages that pay tribute in silence and spies.
\end{itemize}

\section*{28. Nihori — The Isles of the Dawn Spirit}

\subsection*{Overview}
A storm‑wracked archipelago of sea cliffs, bamboo forests, and fire‑mountains. Nihori is ruled by a divine emperor, but his will is carried out by shōguns, daimyōs, and warriors who live by the code of the blade.

\subsection*{Key Features}
\begin{itemize}[leftmargin=*]
    \item The \textbf{Inland Sea} – a calm heart of the isles, dotted with shrine‑islands and pirate nests.
    \item The \textbf{Fire‑Mountains} – volcanoes that smoke year‑round; their ash enriches the soil and curses the unlucky.
    \item The Shōgun – a military dictator who rules in the emperor’s name; his palace is a floating fortress of black lacquer.
    \item The Samurai – a warrior caste bound by honour and debt; their swords are registered like deeds.
    \item The Shinobi – spies and assassins who serve no clan, only the highest bidder; their tools are poetry and poison.
\end{itemize}

\section*{29. Ayokha — The Monsoon Throne, The River of Heaven}

\subsection*{Overview}
A sprawling jungle kingdom ruled by a god‑king who is also the monsoon. Ayokha is a land of river barges, elephant cavalry, and temples that flood and drain with the seasons.

\subsection*{Key Features}
\begin{itemize}[leftmargin=*]
    \item The \textbf{Sona River} – a slow, hieroglyphic waterway that changes course each year; maps are valid only for the season.
    \item The \textbf{Jade Coast} – a strip of emerald beaches, where traders from Sihai and Nihori meet to haggle over silk and sapphires.
    \item War Elephants – armoured behemoths whose howdahs are mobile fortresses; their mahouts never speak above a whisper.
    \item The Royal Guard – warrior monks who wear wooden armour and carry no metal; they believe steel attracts storms.
    \item Monsoon‑riding junks – ships with retractable masts, designed to surf the floodwaters inland.
\end{itemize}

\section*{30. Alberriden Sea — Cold Mirror of the North}

\subsection*{Overview}
A brackish, iron‑skied sea that mirrors the north. The Alberriden is a place of fog, mountain winds, and root‑cellars carved into cliffs.

\subsection*{Key Features}
\begin{itemize}[leftmargin=*]
    \item The \textbf{Yrolka Mouth} – the great river delta where Linn longships beach to trade furs for grain.
    \item The \textbf{Brack Marsh Rim} – a salt marsh that hosts wading birds and will‑o’‑wisps.
    \item Valewood Edge – a coast of silver‐barked trees that lean over the water, their roots drinking salt.
    \item Mistland Coast – where the dwarven bell‑lines meet the sea; the final levee is a lighthouse called the Quiet Bell.
    \item The Haravoa Shadow – a deep trench where the water is black and the fish have no eyes.
\end{itemize}

\section*{31. The Crimson Basin}

\subsection*{Overview}
A vast rainforest contested between wood elves and Oshiiran settlers. The Crimson Basin is a realm of rivers, treaties, and farm‑forests that bleed red sap when cut.

\subsection*{Key Features}
\begin{itemize}[leftmargin=*]
    \item The \textbf{Enjwe Trunk} – a colossal silk‑cotton tree whose canopy covers a square mile; its roots are a market.
    \item The \textbf{Heartwood} – a sacred grove where the Lethai‑ar bury their dead; the ground is quilted with moss and memory.
    \item Wood Elf longhouses raised on stilts; their walls are woven from living branches.
    \item Oshiiran basin‑folk who grow coffee and indigo on terraced slopes, paying taxes in beans and secrecy.
    \item The \textbf{Treaty of Three Waters} – a pact that divides the basin’s rivers into zones of influence; it is broken every decade.
\end{itemize}

\section*{32. Linnstad}

\subsection*{Overview}
A northern city‑state and fur‑trade hub, the southernmost outpost of the Linnic tribes. Linnstad claims to have invented the dragon boat – a claim disputed by every other port on the Yloka.

\subsection*{Key Features}
\begin{itemize}[leftmargin=*]
    \item A palisade of whale ribs, tarred black; inside, longhouses share walls with counting‑houses.
    \item The fur market: ermine, sable, and the occasional white fox pelt – each sold with a story of the hunt.
    \item The Yrolka Mouth – a confluence of rivers where Linnstad’s boatwrights cut the best keels.
    \item Fierce, pragmatic culture: oaths are witnessed by a notary and a skald; both must agree for the vow to bind.
\end{itemize}

\section*{33. Rabelle}

\subsection*{Overview}
A mountain city of ore and gem wealth, ruled by red‑haired Rabellans who claim giant ancestry. Rabelle’s forges never sleep, and its passes are treacherous.

\subsection*{Key Features}
\begin{itemize}[leftmargin=*]
    \item Fine metalwork – the Rabellans weave iron as if it were cloth, producing mail so light it is mistaken for silk.
    \item Dwarven trade: Rabelle is the only human city permitted to mint coins with Aeler runes.
    \item The pass toward the Mistlands is a knife‑edge ridge; guides charge triple during fog season.
    \item Blunt pride: a Rabelle merchant will tell you exactly what your goods are worth, then spit at your feet. This is considered good manners.
\end{itemize}

\section*{34. Northpass}

\subsection*{Overview}
A frontier town and gateway to the Mistlands. Northpass is a bare‑bones outpost of vigilance – a place where you stock up on salt and silence before crossing the levee.

\subsection*{Key Features}
\begin{itemize}[leftmargin=*]
    \item The \textbf{Ermine Inn} – a timber‑and‑thatch hostel where every guest is registered by bell‑code.
    \item Caravan stops: wagons pay a toll in iron nails (used to repair the levee bells).
    \item Carefully monitored by dwarves: a stone bunkhouse houses a garrison of Spirit Shield Warriors.
    \item Rough edges: the watch commander has hanged three smugglers this month; he keeps their boots as planters.
\end{itemize}

\section*{35. The Northern \& Eastern City-States}

\subsection*{Overview}
Independent city‑states, mountain strongholds, and coastal havens that answer to no crown. They are bound by trade, survival, and the mutual need for neutrality.

\subsection*{Key Features}
\begin{itemize}[leftmargin=*]
    \item Each state has its own coin, its own law, and its own god – sometimes three of each.
    \item Dangerous, decadent, or dishonourable – but everyone trades. A letter of credit from one bank is honoured in a rival’s counting‑house.
    \item Known for the dangers and rewards of the frontier: monster‑hunting, relic‑grabbing, and the occasional assassination disguised as a business dispute.
    \item The League of Freeholds – a loose alliance that meets once a year on a bridge between two states; the bridge has no national flag, only a toll booth.
\end{itemize}

\end{document}